\DocumentMetadata{
  pdfversion=2.0,
  pdfstandard=ua-2,
  lang=en-US,
  tagging=on,
  tagging-setup={math/setup=mathml-SE,tabsorder=structure}
}
\documentclass[11pt,letterpaper]{article}
\usepackage[margin=0.68in,includefoot]{geometry}
\usepackage{newcomputermodern}
\usepackage{amsmath,amscd,mathtools}
\usepackage{tikz,adjustbox}
\providecommand{\square}{\mdlgwhtsquare}
\usepackage[hidelinks]{hyperref}
\hypersetup{
  pdftitle={Numerical Linear Algebra First-Year Exam, August 2024},
  pdfauthor={University of Florida Department of Mathematics},
  pdfsubject={Graduate mathematics examination},
  pdfdisplaydoctitle=true,
  bookmarksopen=true
}
\setcounter{secnumdepth}{0}
\setlength{\parindent}{0pt}
\setlength{\parskip}{0.28em}
\setlength{\emergencystretch}{3em}
\setlength{\leftmargini}{2.15em}
\setlength{\leftmarginii}{2.65em}
\setlength{\leftmarginiii}{3em}
\pagestyle{empty}
\raggedbottom
\allowdisplaybreaks
\definecolor{ProblemConcern}{HTML}{7A1F1F}
\newenvironment{problemconcern}{%
  \par\tagstructbegin{tag=Aside}\begingroup\color{ProblemConcern}
}{%
  \par\endgroup\tagstructend\nopagebreak[4]\vspace{0.12em}
}
\NewTaggingSocketPlug{block/list/label}{examproblem}{%
  \tagstructbegin{tag=Lbl}\tagstructend%
  \tagstructbegin{tag=\LBody}%
  \tagstructbegin{tag=\ProblemHeadingTag,title-o={\ProblemHeadingText}}%
  \tagmcbegin{tag=\ProblemHeadingTag,actualtext={\ProblemHeadingText}}%
  #2%
  \tagmcend%
  \tagstructend%
}
\newenvironment{examproblems}[2]{%
  \def\ProblemHeadingTag{#2}%
  \AssignTaggingSocketPlug{block/list/label}{examproblem}%
  \begin{enumerate}%
  \setcounter{enumi}{#1}%
  \renewcommand{\labelenumi}{\textbf{\theenumi.}}%
  \setlength{\topsep}{0.35em}%
  \setlength{\partopsep}{0pt}%
  \setlength{\itemsep}{0.48em}%
  \setlength{\parsep}{0.18em}%
}{\end{enumerate}}
\newenvironment{examsubparts}{%
  \AssignTaggingSocketPlug{block/list/label}{default}%
  \begin{enumerate}%
  \setlength{\topsep}{0.18em}%
  \setlength{\partopsep}{0pt}%
  \setlength{\itemsep}{0.16em}%
  \setlength{\parsep}{0.1em}%
}{\end{enumerate}}
\newenvironment{examinstructions}{%
  \par\begingroup\itshape
}{%
  \par\endgroup\vspace{0.18em}
}
\makeatletter
\renewcommand\subsection{\@startsection{subsection}{2}{0pt}%
  {-1.25ex plus -.25ex minus -.1ex}%
  {0.55ex plus .1ex}%
  {\normalfont\large\bfseries}}
\renewcommand{\@maketitle}{%
  \newpage\null\vskip -1em%
  {\footnotesize\bfseries
    \href{https://www.ufl.edu/}{University of Florida}, \href{https://math.ufl.edu/}{Department of Mathematics}\par}%
  \vskip -0.5em%
  \tagstructbegin{tag=H1,title={Numerical Linear Algebra First-Year Exam, August 2024}}%
  \tagpdfparaOff%
  \tagmcbegin{tag=H1}%
    {\LARGE\bfseries\href{https://gma.math.ufl.edu/exams/numerical-linear-algebra-first-year/}{Numerical Linear Algebra First-Year Exam}, August 2024\par}%
  \tagmcend%
  \tagpdfparaOn%
  \tagstructend%
  \vskip 0.25em%
  \tagmcbegin{artifact}%
    \hrule height 1pt%
  \tagmcend%
  \vskip 1em%
}
\makeatother
\title{Numerical Linear Algebra First-Year Exam, August 2024}
\author{}
\date{}
\begin{document}
\maketitle
\thispagestyle{empty}
\begin{examinstructions}
Do 4 (four) problems.
\end{examinstructions}
\begin{examproblems}{0}{H2}
\def\ProblemHeadingText{Problem 1.}
\item
Let~\(A\) be an \(m\times n\) matrix (\(m\ge n\)) and let \(A=\widehat Q\widehat R\) be a reduced \(QR\) factorization of~\(A\). Prove that~\(A\) has full rank if and only if the diagonal entries of \(\widehat R\) are nonzero.
\def\ProblemHeadingText{Problem 2.}
\item
\begin{problemconcern}
\textbf{Warning: Suspected error.} Part (c) is preserved as printed, but the stated assumptions do not ensure that \(A^{-1}\) exists; for example, \(A=0\) satisfies \(\|A\|_p<1\). The intended additional assumption or alteration is not uniquely determined.
\end{problemconcern}
Assume that \(A\in\mathbb C^{n\times n}\) and there exists \(p\ge 1\) such that \(\|A\|_p<1\), where \(\|\cdot\|_p\) is a vector-induced matrix norm.
\begin{examsubparts}
\item[\textnormal{(a)}]
Prove that \(I-A\) is invertible.
\item[\textnormal{(b)}]
Prove that
\[
(I-A)^{-1}=\sum_{k=0}^{\infty}A^k.
\]
\item[\textnormal{(c)}]
Prove that \(\|A\|_q\|A^{-1}\|_q\ge 1\), \(\forall\,1\le q\le\infty\).
\item[\textnormal{(d)}]
Prove that
\[
\frac{1}{1+\|A\|_p}\le \|(I-A)^{-1}\|_p\le\frac{1}{1-\|A\|_p}.
\]
\end{examsubparts}
\def\ProblemHeadingText{Problem 3.}
\item
Suppose~\(A\) is Hermitian positive definite.
\begin{examsubparts}
\item[\textnormal{(a)}]
Prove that each principal submatrix of~\(A\) is Hermitian positive definite.
\item[\textnormal{(b)}]
Prove that an element of~\(A\) with largest magnitude lies on the diagonal.
\item[\textnormal{(c)}]
Prove that~\(A\) has a Cholesky decomposition.
\end{examsubparts}
\def\ProblemHeadingText{Problem 4.}
\item
Let \(\mathbf v=[2,-1,1]^T\) and \(H=(\operatorname{span}\{\mathbf v\})^\perp\).
\begin{examsubparts}
\item[\textnormal{(a)}]
Find the matrix \(P_{\mathbf v}\), the orthogonal projection onto \(\operatorname{span}\{\mathbf v\}\).
\item[\textnormal{(b)}]
Find the matrix \(P_H\), the orthogonal projection onto~\(H\).
\item[\textnormal{(c)}]
Find \(Q_H\), the unitary matrix that reflects across~\(H\).
\end{examsubparts}
\def\ProblemHeadingText{Problem 5.}
\item
Let \(A,\delta A\in\mathbb R^{n\times n}\) be full rank and \(b,x\) and \(\delta x\in\mathbb R^n\). Prove that if
\[
Ax=b,\quad\text{and}\quad(A+\delta A)(x+\delta x)=b,
\]
 then
\[
\frac{\|\delta x\|}{\|x+\delta x\|}\le\kappa(A)\frac{\|\delta A\|}{\|A\|},
\]
 where \(\kappa(A)\) is the condition number of~\(A\).
\end{examproblems}
\end{document}
